\documentclass{../../note}

\usepackage{xcolor}
\usepackage{longtable}
\usepackage{array}
\usepackage{listings}

\lstset{
  basicstyle=\ttfamily\small,
  keywordstyle=\color{blue},
  commentstyle=\color{green!50!black},
  stringstyle=\color{red!60!black},
  numbers=left,
  numberstyle=\tiny,
  numbersep=6pt,
  breaklines=true,
  frame=single
}

\title{HUFE 数据结构第一讲}
\author{isomoes}

\begin{document}

\maketitle

\section{本次课定位}

本次课对应 5 次课安排中的第 1 次课，时长 3 小时，目标是完成\textbf{数据结构绪论、线性表、栈和队列}的基础梳理，并通过典型题建立做题入口。

\subsection{本次课目标}

\begin{itemize}
  \item 理解数据、数据元素、数据结构、数据类型、抽象数据类型等基础概念。
  \item 区分逻辑结构与存储结构，建立“逻辑关系 + 存储实现 + 基本操作”的整体视角。
  \item 掌握线性表的定义，理解顺序表与链表的差异及常见考法。
  \item 掌握栈和队列的结构特征、基本操作、顺序与链式实现。
  \item 能处理基础复杂度判断、结构选择、循环队列判空判满、栈队列应用题。
\end{itemize}

\subsection{考试导向提醒}

这一讲虽然属于入门内容，但在试卷中出现频率很高，常见出题方式有三类：

\begin{itemize}
  \item \textbf{概念判断题}：考查逻辑结构、存储结构、ADT、复杂度等基本概念。
  \item \textbf{比较分析题}：比较顺序表与链表、栈与队列、顺序存储与链式存储的优缺点。
  \item \textbf{基础计算题}：如顺序表插删复杂度、循环队列长度、栈的合法出栈序列判断。
\end{itemize}

因此复习时不要只背定义，还要形成“\textbf{定义 - 特点 - 操作 - 应用 - 易错点}”的完整链条。

\subsection{3小时建议节奏}

\begin{center}
\begin{tabular}{|c|c|p{9cm}|}
\hline
阶段 & 时间 & 主要内容 \\
\hline
第1段 & 30分钟 & 绪论：基本概念、逻辑结构与存储结构、算法与复杂度入门 \\
\hline
第2段 & 60分钟 & 线性表：定义、ADT、顺序表、链表、典型比较题 \\
\hline
第3段 & 70分钟 & 栈和队列：定义、基本操作、顺序/链式实现、循环队列 \\
\hline
第4段 & 20分钟 & 栈和队列应用：括号匹配、表达式、层次遍历、调度 \\
\hline
\end{tabular}
\end{center}

\section{绪论}

\subsection{核心概念}

\begin{enumerate}
  \item \textbf{数据}：对客观事物的符号表示，是信息的载体。
  \item \textbf{数据元素}：数据的基本单位，一个数据元素通常可由若干数据项组成。
  \item \textbf{数据结构}：相互之间存在特定关系的数据元素的集合。
  \item \textbf{数据类型}：一组性质相同的值及其上的操作的总称。
  \item \textbf{抽象数据类型（ADT）}：从“逻辑功能”角度描述数据对象和操作，而不关心具体实现。
\end{enumerate}

\subsection{数据结构的三个层次}

学习一个数据结构时，建议固定从以下三个角度去看：

\begin{itemize}
  \item \textbf{逻辑结构}：元素之间是什么关系。
  \item \textbf{存储结构}：这种关系在计算机中怎样存。
  \item \textbf{基本操作}：查找、插入、删除、访问如何实现，代价如何。
\end{itemize}

\subsection{逻辑结构分类}

\begin{itemize}
  \item 集合结构：元素之间除“同属一个集合”外无其他关系。
  \item 线性结构：元素之间是一对一关系。
  \item 树形结构：元素之间是一对多关系。
  \item 图形结构：元素之间是多对多关系。
\end{itemize}

\subsection{存储结构分类}

\begin{itemize}
  \item 顺序存储：地址连续，支持随机访问。
  \item 链式存储：地址可不连续，依靠指针联系。
  \item 索引存储：在数据之外附加索引信息。
  \item 散列存储：根据关键字直接计算存储位置。
\end{itemize}

\subsection{算法与复杂度入门}

\textbf{算法}是求解问题的一组有穷、确定、可行的步骤。算法通常关注两个指标：

\begin{itemize}
  \item \textbf{时间复杂度}：执行操作次数随问题规模 $n$ 增长的变化趋势。
  \item \textbf{空间复杂度}：算法运行额外占用的存储空间规模。
\end{itemize}

常见时间复杂度从优到劣可记为：

\[
O(1) < O(\log n) < O(n) < O(n\log n) < O(n^2) < O(n^3)
\]

\subsection{复杂度理解例子}

\begin{enumerate}
  \item 顺序表按位访问第 $i$ 个元素，只需一次地址计算，因此常看作 $O(1)$。
  \item 在长度为 $n$ 的顺序表表头插入一个元素，通常需要移动原有 $n$ 个元素，因此为 $O(n)$。
  \item 在单链表表头插入一个新结点，只需修改少量指针，因此若头指针已知，可看作 $O(1)$。
  \item 在单链表中查找第 $i$ 个结点，需要从头结点开始逐个后移，通常为 $O(n)$。
\end{enumerate}

\subsection{绪论部分常见判断}

\begin{enumerate}
  \item 逻辑结构和存储结构不是一一对应关系；同一逻辑结构可对应不同存储结构。
  \item 抽象数据类型强调“做什么”，具体数据结构强调“怎么做”。
  \item 时间复杂度只关注数量级，不关注常数项和低阶项。
\end{enumerate}

\subsection{绪论练习}

\begin{enumerate}
  \item 什么是数据结构？请用“数据元素关系”的角度作答。
  \item 数据的逻辑结构与存储结构有什么区别？
  \item 为什么说抽象数据类型和具体实现不是同一个层面的问题？
  \item 请判断：线性结构一定采用顺序存储。（\hspace{1cm}）
  \item 请写出下列复杂度从小到大的排列：$O(n)$、$O(1)$、$O(n^2)$、$O(\log n)$。
\end{enumerate}

\section{线性表}

\subsection{定义}

\textcolor{red}{线性表}是具有相同数据类型的 $n$ 个数据元素的有限序列。

其特点是：

\begin{itemize}
  \item 除第一个元素外，每个元素有且仅有一个直接前驱。
  \item 除最后一个元素外，每个元素有且仅有一个直接后继。
  \item 当 $n=0$ 时，称为空表。
\end{itemize}

\subsection{线性表的ADT视角}

从考题角度，线性表的 ADT 常见操作包括：

\begin{itemize}
  \item InitList：初始化
  \item ListEmpty：判空
  \item ListLength：求长度
  \item GetElem：按位查找
  \item LocateElem：按值查找
  \item ListInsert：插入
  \item ListDelete：删除
\end{itemize}

这里要注意两组容易混淆的操作：

\begin{itemize}
  \item \textbf{按位查找}：根据位置找元素，如找第 $i$ 个元素。
  \item \textbf{按值查找}：根据元素值找其所在位置或结点。
\end{itemize}

在顺序表中，按位查找优势明显；在链表中，两者通常都需要遍历。

\subsection{顺序表}

\subsubsection{定义与特点}

顺序表使用一组\textbf{地址连续}的存储单元依次存储线性表元素，因此逻辑上相邻的元素在物理上也相邻。

\textbf{优点：}
\begin{itemize}
  \item 支持随机访问，按位查找快。
  \item 存储密度高，不需要额外指针域。
\end{itemize}

\textbf{缺点：}
\begin{itemize}
  \item 插入和删除通常需要移动大量元素。
  \item 容量预先确定，扩展不灵活。
\end{itemize}

\subsubsection{基本复杂度}

\begin{center}
\begin{tabular}{|c|c|c|}
\hline
操作 & 最好情况 & 最坏情况 \\
\hline
按位查找 & $O(1)$ & $O(1)$ \\
\hline
按值查找 & $O(1)$ & $O(n)$ \\
\hline
插入 & $O(1)$ & $O(n)$ \\
\hline
删除 & $O(1)$ & $O(n)$ \\
\hline
\end{tabular}
\end{center}

\subsubsection{典型考点}

\begin{enumerate}
  \item 在第 $i$ 个位置插入元素时，需要将第 $i$ 个及其后的元素整体后移。
  \item 在表头插入或删除，最坏情况下都需要移动 $n$ 级别元素。
  \item 顺序表适合查询多、修改少的场景。
\end{enumerate}

\subsubsection{顺序表操作过程示意}

设原顺序表为：

\[
(a_1, a_2, a_3, a_4)
\]

若在第 3 个位置插入元素 $x$，插入后得到：

\[
(a_1, a_2, x, a_3, a_4)
\]

此时原来的 $a_3, a_4$ 都要向后移动一个位置。可见，\textbf{顺序表插入删除的核心代价往往来自元素移动，而不是赋值本身}。

\subsection{链表}

\subsubsection{定义与特点}

链表用一组\textbf{地址可以不连续}的存储单元存放元素，通过指针反映逻辑关系。

\textbf{优点：}
\begin{itemize}
  \item 插入、删除时不需要整体移动元素。
  \item 空间可动态申请，扩展方便。
\end{itemize}

\textbf{缺点：}
\begin{itemize}
  \item 不支持随机访问，按位查找通常要从头遍历。
  \item 每个结点需要额外指针域，存储密度较低。
\end{itemize}

\subsubsection{常见类型}

\begin{itemize}
  \item 单链表：每个结点只有一个后继指针。
  \item 双链表：每个结点有前驱和后继两个指针。
  \item 循环链表：尾结点指针不为空，而是重新指向表头或头结点。
\end{itemize}

\subsubsection{顺序表与链表对比}

\begin{center}
\begin{tabular}{|>{\centering\arraybackslash}m{2.6cm}|>{\centering\arraybackslash}m{4cm}|>{\centering\arraybackslash}m{4cm}|}
\hline
比较项目 & 顺序表 & 链表 \\
\hline
存储空间 & 连续 & 不连续 \\
\hline
按位访问 & 快，$O(1)$ & 慢，$O(n)$ \\
\hline
插入删除 & 常需移动元素 & 常只需修改指针 \\
\hline
存储密度 & 高 & 较低 \\
\hline
适用场景 & 查询多 & 插删多 \\
\hline
\end{tabular}
\end{center}

\subsubsection{单链表插入与删除思想}

若已知结点 $p$，要在其后插入新结点 $s$：

\[
s \rightarrow next = p \rightarrow next, \quad p \rightarrow next = s
\]

若删除结点 $p$ 的后继结点 $q$：

\[
p \rightarrow next = q \rightarrow next
\]

注意：单链表删除某结点通常必须先找到其前驱结点。

\subsubsection{链表操作过程理解}

设有单链表：

\[
head \rightarrow a \rightarrow b \rightarrow c \rightarrow NULL
\]

若要在结点 $b$ 后插入结点 $x$，操作顺序不能写反：

\begin{enumerate}
  \item 先让 $x$ 指向 $b$ 的原后继，即 $x \rightarrow next = b \rightarrow next$；
  \item 再让 $b$ 指向 $x$，即 $b \rightarrow next = x$。
\end{enumerate}

若先改第二步，就会丢失原后继结点地址。这类“\textbf{指针修改顺序题}”很容易在选择题和填空题中出现。

\subsubsection{线性表部分的选型思路}

当题目问“应该选择哪种结构”时，可以按下面思路判断：

\begin{itemize}
  \item 若题目强调\textbf{随机访问、按位置快速读取}，优先考虑顺序表。
  \item 若题目强调\textbf{频繁插入删除、长度变化较大}，优先考虑链表。
  \item 若题目强调\textbf{存储密度高}，顺序表通常更优。
  \item 若题目强调\textbf{不要求连续空间}，链表更灵活。
\end{itemize}

\subsection{线性表练习}

\begin{enumerate}
  \item 线性表的两个基本逻辑特征是什么？
  \item 顺序表为什么适合按位查找？
  \item 单链表为什么不适合频繁按位访问？
  \item 在顺序表第 1 个位置插入元素，为什么最坏时间复杂度为 $O(n)$？
  \item 在单链表中，若已知某结点前驱，则插入一个新结点的时间复杂度通常是多少？
  \item 请判断：链表的插入操作一定比顺序表快。（\hspace{1cm}）
\end{enumerate}

\section{栈和队列}

\subsection{栈}

\subsubsection{定义}

栈是限定仅在表尾进行插入和删除的线性表。允许操作的一端称为\textbf{栈顶}，另一端称为\textbf{栈底}。

栈的核心特征是：\textcolor{red}{后进先出（LIFO）}。

\subsubsection{基本操作}

\begin{itemize}
  \item InitStack
  \item StackEmpty
  \item Push
  \item Pop
  \item GetTop
\end{itemize}

\subsubsection{顺序栈}

顺序栈通常使用数组和栈顶指针 \texttt{top} 实现。

\begin{lstlisting}[language=C]
#define MaxSize 100
typedef struct {
    ElemType data[MaxSize];
    int top;
} SqStack;
\end{lstlisting}

常见约定：

\begin{itemize}
  \item 初始化：\texttt{top = -1}
  \item 栈空：\texttt{top == -1}
  \item 栈满：\texttt{top == MaxSize - 1}
  \item 进栈：先 \texttt{top++}，再存值
  \item 出栈：先取值，再 \texttt{top--}
\end{itemize}

如果栈中当前有 3 个元素，则栈顶指针通常指向第 3 个元素的位置，而不是下一个空位置。这一点要和某些教材中的不同约定区分开来。

\subsubsection{链栈}

链栈本质上是操作受限的单链表，通常在链表头部进行入栈和出栈。

\begin{lstlisting}[language=C]
typedef struct LinkNode {
    ElemType data;
    struct LinkNode *next;
} LinkNode, *LinkStack;
\end{lstlisting}

链栈的优点是扩容灵活，不需要预先分配固定长度。

\subsubsection{栈的常见应用}

\begin{itemize}
  \item 括号匹配
  \item 表达式求值
  \item 函数调用与递归
  \item 浏览器前进后退
\end{itemize}

\subsubsection{为什么表达式求值常用栈}

表达式求值之所以常用栈，是因为运算具有“后出现的高优先级运算可能先处理”的特点，和栈的后进先出特征非常契合。

例如中缀表达式：

\[
3 + 4 \times 2
\]

虽然 $+$ 先出现，但计算时必须先处理 $4 \times 2$。因此在实现中，常用一个栈保存运算符，必要时也会再用一个栈保存操作数。

\subsubsection{栈的典型判断方法}

判断一个序列是否可能为合法出栈序列时，最稳妥的方法不是凭感觉，而是\textbf{模拟入栈和出栈过程}：

\begin{enumerate}
  \item 按给定顺序依次入栈；
  \item 每入一个元素，就检查栈顶是否等于目标出栈序列当前元素；
  \item 若相等则持续出栈；
  \item 最终若目标序列可全部匹配，则该序列合法。
\end{enumerate}

\subsection{栈练习}

\begin{enumerate}
  \item 栈为什么称为操作受限的线性表？
  \item 栈的逻辑特征是什么？请用一句话概括。
  \item 顺序栈中，\texttt{top = -1} 一般表示什么状态？
  \item 入栈序列为 $1,2,3$，请写出两个合法的出栈序列。
  \item 为什么括号匹配问题通常使用栈而不是队列？
  \item 请判断：函数递归调用过程中不需要栈空间。（\hspace{1cm}）
\end{enumerate}

\subsection{队列}

\subsubsection{定义}

队列是只允许在一端插入、另一端删除的线性表。插入的一端称为\textbf{队尾}，删除的一端称为\textbf{队头}。

队列的核心特征是：\textcolor{red}{先进先出（FIFO）}。

\subsubsection{基本操作}

\begin{itemize}
  \item InitQueue
  \item QueueEmpty
  \item EnQueue
  \item DeQueue
  \item GetHead
\end{itemize}

\subsubsection{循环队列是重点}

普通顺序队列会出现“假溢出”问题，因此考试中更常考\textbf{循环队列}。

\begin{lstlisting}[language=C]
#define MaxSize 100
typedef struct {
    ElemType data[MaxSize];
    int front, rear;
} SqQueue;
\end{lstlisting}

采用“牺牲一个存储单元”方案时：

\begin{itemize}
  \item 初始化：\texttt{front = rear = 0}
  \item 队空：\texttt{front == rear}
  \item 队满：\texttt{(rear + 1) \% MaxSize == front}
  \item 入队：先存入 \texttt{data[rear]}，再更新 \texttt{rear}
  \item 出队：先取 \texttt{data[front]}，再更新 \texttt{front}
\end{itemize}

这里的 \texttt{front} 一般指向当前队头元素的位置，\texttt{rear} 一般指向下一个可插入位置。正因为采用了这种定义，才需要牺牲一个存储单元来区分“空”和“满”。

队列长度公式：

\[
\text{length} = (rear - front + MaxSize) \% MaxSize
\]

\subsubsection{循环队列计算示例}

若 \texttt{MaxSize = 6}，当前 \texttt{front = 4}，\texttt{rear = 1}，则说明队列发生了“回绕”。此时长度为：

\[
(1 - 4 + 6) \% 6 = 3
\]

这说明循环队列的长度计算不能简单用 \texttt{rear - front}，必须使用统一公式。

\subsubsection{链队}

链队通常带头结点，\texttt{front} 指向队头，\texttt{rear} 指向队尾。

\begin{itemize}
  \item 入队在队尾进行。
  \item 出队在队头进行。
  \item 若最后一个元素出队，需让 \texttt{rear = front}。
\end{itemize}

\subsubsection{队列的常见应用}

\begin{itemize}
  \item 层次遍历
  \item 广度优先搜索
  \item 作业调度
  \item 打印任务管理
  \item 消息缓冲
\end{itemize}

\subsubsection{为什么层次遍历要用队列}

二叉树层次遍历要求“先访问上一层，再访问下一层；同层中先左后右”。这正好符合队列的先进先出特征：

\begin{enumerate}
  \item 根结点先入队；
  \item 先出队的结点先访问；
  \item 再把它的左右孩子按顺序入队；
  \item 这样就能保证结点按层次顺序被访问。
\end{enumerate}

\subsection{队列练习}

\begin{enumerate}
  \item 队列为什么也是操作受限的线性表？
  \item 队列的逻辑特征是什么？
  \item 普通顺序队列为什么会出现“假溢出”？
  \item 在循环队列中，采用牺牲一个存储单元的方法时，队空和队满条件分别是什么？
  \item 若 \texttt{MaxSize = 8}，\texttt{front = 6}，\texttt{rear = 2}，则循环队列长度是多少？
  \item 请判断：广度优先搜索通常使用队列实现。（\hspace{1cm}）
\end{enumerate}

\section{易混点总结}

\begin{enumerate}
  \item \textbf{线性表、栈、队列的关系}：栈和队列都是特殊的线性表。
  \item \textbf{顺序表和数组的关系}：顺序表通常借助数组实现，但顺序表强调的是一种线性表存储方式。
  \item \textbf{链表按位访问慢}：因为逻辑位置无法直接映射为物理地址。
  \item \textbf{循环队列判满}：若采用牺牲一个单元策略，不能把 \texttt{front == rear} 同时看作“空”和“满”。
  \item \textbf{栈与递归}：递归的本质通常依赖栈保存现场。
\end{enumerate}

\section{本讲知识框架}

建议把本讲内容整理成下面这条主线：

\begin{center}
\begin{tabular}{|c|c|c|c|}
\hline
主题 & 核心问题 & 典型结构 & 高频考点 \\
\hline
绪论 & 什么是数据结构 & 逻辑结构/存储结构 & 概念辨析、复杂度 \\
\hline
线性表 & 线性关系如何存储 & 顺序表、链表 & 对比、选型、插删代价 \\
\hline
栈 & 为什么后进先出 & 顺序栈、链栈 & 出栈序列、括号匹配 \\
\hline
队列 & 为什么先进先出 & 循环队列、链队 & 判空判满、长度计算 \\
\hline
\end{tabular}
\end{center}

\section{典型例题}

\subsection{例题1：顺序表插入的复杂度}

\textbf{题目：}在长度为 $n$ 的顺序表第 1 个位置插入一个新元素，其时间复杂度是多少？

\textbf{分析：}原有元素都要后移，共涉及 $n$ 级别移动。

\textbf{答案：}$O(n)$。

\subsection{例题2：链表与顺序表选择}

\textbf{题目：}若某应用系统需要频繁插入和删除元素，而较少按位随机访问，应优先使用顺序表还是链表？

\textbf{分析：}插删频繁时，链表不需要大量移动元素，更合适。

\textbf{答案：}链表。

\subsection{例题3：栈的出栈序列判断}

\textbf{题目：}入栈序列为 $1,2,3,4$，下列序列中哪一个不可能是出栈序列？

备选：
\begin{itemize}
  \item A. $4,3,2,1$
  \item B. $3,4,2,1$
  \item C. $2,4,3,1$
  \item D. $3,1,4,2$
\end{itemize}

\textbf{分析：}D 不可能。因为 3 出栈后，1 若要先于 4 出栈，则 4 必须尚未入栈；但 2 已在 1 前入栈，之后再让 4 在 2 后出栈会矛盾。

\textbf{答案：}D。

\subsection{例题4：循环队列长度}

\textbf{题目：}循环队列中 \texttt{MaxSize = 8}，当前 \texttt{front = 2}，\texttt{rear = 6}，则队列中元素个数为多少？

\textbf{解：}
\[
(6 - 2 + 8) \% 8 = 4
\]

\textbf{答案：}4 个。

\subsection{例题5：链表删除的关键条件}

\textbf{题目：}在单链表中删除结点 $p$ 的后继结点，操作的关键前提是什么？

\textbf{分析：}必须保证结点 $p$ 存在且其后继结点非空，否则无法完成删除。

\textbf{答案：}已知待删结点的前驱结点，且前驱结点存在后继。

\subsection{例题6：结构选择题}

\textbf{题目：}某系统需要频繁读取第 $i$ 个元素，但很少做插入和删除，应选链表还是顺序表？

\textbf{分析：}频繁按位访问时，顺序表可直接通过地址计算定位。

\textbf{答案：}顺序表。

\section{课堂练习}

\subsection{基础练习}

\begin{enumerate}
  \item 数据结构从通常意义上看包括哪三个方面？
  \item 逻辑结构可分为哪四类？
  \item 顺序表按位查找为什么是 $O(1)$？
  \item 单链表删除某结点时，为什么通常需要先找到其前驱？
  \item 栈和队列分别满足什么原则？
\end{enumerate}

\subsection{判断与选择}

\begin{enumerate}
  \item 同一种逻辑结构只能对应一种存储结构。（\hspace{1cm}）
  \item 链表的插入和删除一定都是 $O(1)$。（\hspace{1cm}）
  \item 栈属于操作受限的线性表。（\hspace{1cm}）
  \item 若循环队列采用牺牲一个存储单元的方法，则 \texttt{front == rear} 表示队满。（\hspace{1cm}）
  \item 队列适合处理具有“先到先服务”特征的问题。（\hspace{1cm}）
  \item 栈顶指针 \texttt{top} 总是指向下一个可用空位置。（\hspace{1cm}）
\end{enumerate}

\subsection{计算与分析}

\begin{enumerate}
  \item 一个顺序表长度为 $n$，在表尾插入元素的最好时间复杂度和最坏时间复杂度分别是多少？
  \item 设循环队列 \texttt{MaxSize = 10}，\texttt{front = 7}，\texttt{rear = 2}，求当前队列长度。
  \item 入栈序列为 $a,b,c,d,e$，请写出一个合法的出栈序列，并说明理由。
  \item 简述顺序表与链表各自更适合什么应用场景。
  \item 为什么说栈和队列都是特殊线性表，但它们的应用场景明显不同？
  \item 若单链表中仅知道某结点本身，而不知道其前驱，删除操作会遇到什么困难？
\end{enumerate}

\subsection{综合小题}

\textbf{题目：}某教务系统需要处理两个功能：

\begin{itemize}
  \item 功能 A：按学号快速读取当前第 $i$ 个学生记录；
  \item 功能 B：模拟打印任务，先提交的任务先打印。
\end{itemize}

请分别为功能 A 和功能 B 选择更合适的数据结构，并说明理由。

\section{课后小测参考答案}

\begin{enumerate}
  \item 数据的逻辑结构、存储结构、基本操作。
  \item 集合结构、线性结构、树形结构、图形结构。
  \item 因为元素地址连续，可由首地址和位序直接计算存储地址。
  \item 因为单链表只有后继指针，没有直接指向前驱的指针。
  \item 栈是后进先出，队列是先进先出。
  \item 第1题错；第2题错；第3题对；第4题错；第5题对；第6题错。
  \item 当 \texttt{MaxSize = 10}，长度为 $(2 - 7 + 10) \% 10 = 5$。
  \item 栈强调“后进入先处理”，适合递归、回退、括号匹配；队列强调“先进入先处理”，适合调度、缓冲、层次遍历。
  \item 因为单链表没有前驱指针，无法直接修改前驱结点的 \texttt{next} 域。
\end{enumerate}

综合小题参考：功能 A 选顺序表，因为按位读取快；功能 B 选队列，因为打印任务符合先进先出原则。

\subsection{分节练习参考答案}

\textbf{绪论练习参考：}

\begin{enumerate}
  \item 数据结构是相互之间存在特定关系的数据元素的集合。
  \item 逻辑结构强调元素之间的逻辑关系；存储结构强调这些关系在计算机中的具体表示。
  \item 因为 ADT 关注“提供哪些操作”，而具体实现关注“这些操作如何落地”。
  \item 错。线性结构也可以采用链式存储。
  \item $O(1) < O(\log n) < O(n) < O(n^2)$。
\end{enumerate}

\textbf{线性表练习参考：}

\begin{enumerate}
  \item 除首元素外每个元素有唯一前驱，除尾元素外每个元素有唯一后继。
  \item 因为顺序表地址连续，可由首地址和位序直接计算存储位置。
  \item 因为链表结点地址不连续，不能直接按位置定位，只能顺序遍历。
  \item 因为原有元素通常都要整体后移。
  \item 若前驱已知，通常为 $O(1)$。
  \item 错。若链表插入前还要先查找位置，则整体代价未必更低。
\end{enumerate}

\textbf{栈练习参考：}

\begin{enumerate}
  \item 因为插入和删除只能在同一端进行。
  \item 后进先出。
  \item 空栈。
  \item 如 $3,2,1$；$2,3,1$。
  \item 因为最近遇到的左括号要最先和后续右括号匹配，符合后进先出。
  \item 错。递归调用通常依赖栈保存返回地址和现场信息。
\end{enumerate}

\textbf{队列练习参考：}

\begin{enumerate}
  \item 因为插入和删除分别限制在两端进行。
  \item 先进先出。
  \item 因为多次出队后，数组前部虽有空位，但 \texttt{rear} 可能已到数组末端。
  \item 队空：\texttt{front == rear}；队满：\texttt{(rear + 1) \% MaxSize == front}。
  \item 长度为 $(2 - 6 + 8) \% 8 = 4$。
  \item 对。
\end{enumerate}

\section{本讲小结}

本讲是数据结构部分的入口，核心不是记定义，而是建立一个稳定的分析框架：\textbf{先看逻辑关系，再看存储实现，最后看基本操作和复杂度}。在线性表部分，要重点比较顺序表和链表；在栈和队列部分，要重点掌握其受限操作特征以及循环队列的判空、判满和长度计算。

下一讲将进入数组和广义表、树和二叉树，难度会明显上升，尤其要提前适应“性质 + 遍历 + 应用题”的出题方式。

\end{document}
